\begin{lemma}
\label{lemma-characterize-local-ring}
Let $R$ be a ring. The following are equivalent:
\begin{enumerate}
\item $R$ is a local ring,
\item $\Spec(R)$ has exactly one closed point,
\item $R$ has a maximal ideal $\mathfrak m$
and every element of $R \setminus \mathfrak m$
is a unit, and
\item $R$ is not the zero ring and for every $x \in R$ either $x$
or $1 - x$ is invertible or both.
\end{enumerate}
\end{lemma}

\begin{proof}
Let $R$ be a ring, and $\mathfrak m$ a maximal ideal.
If $x \in R \setminus \mathfrak m$, and $x$ is not a unit
then there is a maximal ideal $\mathfrak m'$ containing $x$.
Hence $R$ has at least two maximal ideals. Conversely,
if $\mathfrak m'$ is another maximal ideal, then choose
$x \in \mathfrak m'$, $x \not \in \mathfrak m$. Clearly
$x$ is not a unit. This proves the equivalence of (1) and (3).
The equivalence of (1) and (2) follows because the closure of $\{\mathfrak p\}$ is
$V(\mathfrak p)$. Indeed, this is a closed set containing
$\mathfrak p$, and any closed set $V(T)$ containing $\mathfrak p$
has $T\subset\mathfrak p$, hence contains $V(\mathfrak p)$.
A maximal ideal therefore gives a closed point. Conversely, if
$\{\mathfrak p\}$ is closed, apply (2) of
Lemma \ref{lemma-Zariski-topology} to $R/\mathfrak p$.
The inverse image of a maximal ideal of that quotient is a maximal
ideal $\mathfrak m$ of $R$ containing $\mathfrak p$.
Then $\mathfrak m\in V(\mathfrak p)=\{\mathfrak p\}$,
so $\mathfrak m=\mathfrak p$.
If $R$ is local then (4) holds since $x$ is either in $\mathfrak m$
or not. If (4) holds, and $\mathfrak m$, $\mathfrak m'$ are distinct
maximal ideals then we may choose $x \in R$ such that
$x \bmod \mathfrak m' = 0$ and $x \bmod \mathfrak m = 1$
by the Chinese remainder theorem
(Lemma \ref{lemma-chinese-remainder}).
This element $x$ is not invertible and neither is $1 - x$ which is
a contradiction. Thus (4) and (1) are equivalent.
\end{proof}